% A6 + B2-freeze：组件归因
\begin{table}[t]
\centering
\caption{组件级归因：两种互补干预在同一份同分布文本上的结果。
（a）A6 冻结（16\,KB，seed $\in\{0,1\}$）：基线为全部正常写入；
（b）B2-freeze（96\,KB，单条流）在自然重复任务上的 $<\!1$k 桶复现增益。
两组的结论一致：代价几乎全部落在 \texttt{recurrent} 上，矩阵记忆与队列 $+$ MTM 无损失。}
\label{tab:attr}
\small
\begin{tabular}{@{}lcc@{}}
\toprule
\multicolumn{3}{@{}l}{\textbf{(a) A6 冻结单分量（16\,KB，2 seed）}} \\
\midrule
条件 & bits/byte（均值） & $\Delta$ vs 全部写入 \\
\midrule
write（全部写入，基线） & $1.386$ & $0$ \\
lock（冻结全部写入） & $2.420$ & $+1.034$ \\
freeze \texttt{recurrent} & $2.408$ & $\mathbf{+1.022}$ \\
freeze \texttt{memory} & $1.386$ & $-0.0005$ \\
freeze \texttt{queue}$+$\texttt{MTM} & $1.386$ & $+0.0002$ \\
\midrule
\multicolumn{3}{@{}l}{\textbf{(b) B2-freeze：自然重复的 $<\!1$k 桶复现增益（nats，96\,KB）}} \\
\midrule
条件 & bits/byte & $<\!1$k 桶增益 \\
\midrule
write & $1.0838$ & $+0.0839$ \\
lock（全程只读） & $2.1641$ & $-0.0050$ \\
freeze \texttt{recurrent} & $2.1527$ & $\mathbf{-0.0043}$ \\
freeze \texttt{memory} & $1.0836$ & $+0.0840$ \\
freeze \texttt{queue}$+$\texttt{MTM} & $1.0832$ & $+0.0837$ \\
\bottomrule
\end{tabular}
\end{table}
